\chapter{Att kontrollera ett påstående}
\label{app:verifiera}

Det här materialet gör, som allt undervisningsmaterial, en rad
påståenden.  De flesta går att pröva vid datorn: att
\mintinline{python}{1 + frac_a} inte fungerar utan
\mintinline{python}{__radd__} ser ni själva genom att köra programmet.
Några handlar om vad Python \emph{gör} när en operator används, och för
dem är språkets egen dokumentation den rätta källan: de står i
fotnoterna till \cref{sec:addition-subtraktion,sec:forkortning}, med
adress och datum, så att den som vill kan slå upp dem.  Och några
handlar om vad forskningen säger, och dem kan ingen körning avgöra:
\begin{enumerate}
  \item att det är bra att inte upprepa sig i ett program --- principen
    \emph{DRY}, som materialet åberopar när
    \mintinline{python}{__radd__} återanvänder
    \mintinline{python}{__add__} --- och att den kommer från Hunt och
    Thomas (\cref{sec:addition});
  \item att det som varierar mot en invariant bakgrund är det som kan
    urskiljas --- den variationsteoretiska premiss som lärarnoterna i
    marginalen bygger på.
\end{enumerate}
Hur vet man om sådant stämmer?  Samma fråga möter er varje gång ni läser
en lärobok, en webbsida eller ett svar från en språkmodell --- och varje
gång ni själva ska skriva en rapport.  Den här bilagan sammanfattar
arbetsgången.  Båda påståendena ovan är redan prövade i andra
föreläsningar, och \cref{tab:redan-belagt} anger var; arbetsgången
beskrivs utförligare, och genomförd på flera påståenden, i bilagorna
till föreläsningarna \emph{Algoritmiskt tänkande} och \emph{Klasser och
objekt}.

\section{Gör påståendet till en fråga}

Ett påstående går att kontrollera först när det är tydligt vad som skulle
kunna visa att det är \emph{fel}.  Formulera det därför som en fråga med
ett svar.  \enquote{Kommer DRY från Hunt och Thomas?} kan besvaras med
ja eller nej; \enquote{DRY är viktigt} kan det inte, eftersom
\enquote{viktigt} är ett omdöme och inte ett faktum.  Ofta döljer sig
flera påståenden i ett: \enquote{en bra princip som kommer från Hunt
och Thomas} är två frågor --- \emph{ursprung} och \emph{värde} --- och
de kräver olika slags källor.

\section{Välj rätt sorts källa}

Olika frågor besvaras av olika källor.
\begin{description}
  \item[Uppslagsverk] för vad ett ord \emph{betyder}.  För vad ett
    programspråk \emph{gör} är språkets egen dokumentation motsvarande
    källa: att Python provar \mintinline{python}{__radd__} när den
    vänstra operanden inte kan addera står i språkreferensen, inte i
    forskningen.
  \item[Primärkällan] för var en idé \emph{kommer ifrån}, och för ett
    specifikt resultat.
  \item[Översikter] (\foreignlanguage{english}{reviews}) för om något är
    \emph{etablerat}, eller för vad forskningen \emph{sammantaget}
    säger.
  \item[Enskilda studier] för ett specifikt resultat --- med förbehållet
    att en enskild studie bara visar att något observerats en gång, i en
    viss population.
\end{description}
Var man hittar dem: universitetsbibliotekets söktjänst, och databaser som
Scopus, Web of Science, IEEE Xplore, DBLP (datalogi) och OpenAlex
(öppen).

\section{Sök så att du kan ha fel}

Den som söker på namn hon redan känner till hittar det hon väntar sig.
Sök därför \emph{ämnesdrivet} --- på begrepp, inte författare.  Några
praktiska regler:
\begin{itemize}
  \item Håll frågorna korta.  Flera korta frågor slår en lång.
  \item Ställ \emph{samma} frågor i alla databaser.
  \item Sök också efter \emph{invändningar}: lägg till ord som
    \enquote{critique}, \enquote{limitations} eller \enquote{harmful}.
    Ett påstående som överlevt en motsökning är värt mer än ett som
    aldrig prövats, och en invändning som hittas blir ett
    \emph{förbehåll} i svaret, inte ett nederlag.
  \item Räkna med att databasernas rankning missar somligt.  Då hämtas
    källan som \emph{känt dokument}: man vet att den finns och slår upp
    den direkt.  Det är i sin ordning, bara det sägs.
\end{itemize}

\section{Läs i fulltext, och skriv ner hur}

En träff i en databas är inte ett belägg.  Sammanfattningen
(\foreignlanguage{english}{abstract}) räcker inte heller: den säger vad
författarna vill ha sagt, inte vad de visat.  Läs hela texten och leta
upp \emph{det ställe} som styrker påståendet --- ett ordagrant citat med
sidnummer.  Räkna aldrig en källa ni inte läst.

Skriv sedan ner hur ni gick till väga, så att någon annan kan följa
spåret.  I det här materialet står den anteckningen som en kommentar
ovanför varje post i litteraturlistans källfil, med fasta fält:
\begin{description}
  \item[\texttt{CLAIM}] vilket påstående källan ska styrka;
  \item[\texttt{FOUND-VIA}] hur den hittades --- sökfråga och databas,
    eller \enquote{känt dokument};
  \item[\texttt{PICKED}] varför just den, framför andra träffar;
  \item[\texttt{QUOTE}] det ordagranna citatet, med sida;
  \item[\texttt{VERIFIED}] att fulltexten lästs, och var;
  \item[\texttt{COUNTER}] motsökningen och vad den gav;
  \item[\texttt{DATE}] när.
\end{description}
Det tar ett par minuter per källa, och det är det enda som gör det
möjligt att om ett år svara på frågan \enquote{hur vet du det?}

\section{Dra slutsatsen --- med förbehåll}

Svaret skrivs ut tillsammans med det som talar emot.  Att inte upprepa
sig är, enligt bilagan i \emph{Algoritmiskt tänkande}, en risk man tar
på sig medvetet snarare än en regel utan undantag; det är därför
materialet säger att \mintinline{python}{__radd__} \emph{kan} återanvända
\mintinline{python}{__add__} när operationen är kommutativ --- och visar
med \mintinline{python}{__rsub__} att det inte alltid går.  Förbehållen
är inte svagheter i svaret.  De \emph{är} svaret.

\begin{table}[htbp]
  \centering
  \caption{Materialets påståenden om forskningen som frågor, och var
    svaren finns.}
  \label{tab:redan-belagt}
  \begin{tabular}{@{}p{0.5\linewidth}p{0.4\linewidth}@{}}
    \toprule
    Fråga & Svar \\
    \midrule
    Är det bra att inte upprepa sig, och kommer DRY från Hunt och
      Thomas? & Redan belagt i \emph{Algoritmiskt tänkande}, bilaga~E \\
    Är det som varierar mot en invariant bakgrund det som kan
      urskiljas? & Premiss från variationsteorin, med källor i
      lärarnoternas källfil; samma i alla kursens föreläsningar \\
    \bottomrule
  \end{tabular}
\end{table}

\section{En anmärkning om språkmodeller}

Där sökningar gjorts i kursens andra föreläsningar användes en
språkmodell för att \emph{sortera} hundratals träffar --- hör träffen
till ämnet, och åt vilket håll pekar den? --- så att gallringen kunde
granskas.  Modellen sorterade, men varje källa som citeras ska läsas och
väljas av en människa, och även sorteringen ska bekräftas av en
människa.  Använd gärna språkmodeller för att hitta och sortera; låt dem
aldrig vara det sista ledet.
